\documentclass[11pt]{article}

\usepackage[margin=1in]{geometry}
\usepackage{amsmath,amssymb,amsfonts}
\usepackage[hidelinks]{hyperref}
\usepackage{microtype}

\title{Symmetry $d$-matrices: orthogonal, augmented, non-orthogonal}
\author{CoQui developer notes}
\date{2026-05-02}

\newcommand{\bra}[1]{\langle #1 |}
\newcommand{\ket}[1]{| #1 \rangle}
\newcommand{\matel}[3]{\langle #1 | #2 | #3 \rangle}
\newcommand{\AEt}{\mathrm{AE}}
\newcommand{\rotpsiAE}[2]{\hat{R}\,\psi_{#1}^{\,#2,\AEt}}
\newcommand{\psiAE}[2]{\psi_{#1}^{\,#2,\AEt}}
\newcommand{\rotpsipt}[2]{\hat{R}\,\widetilde{\psi}_{#1}^{\,#2}}
\newcommand{\psipt}[2]{\widetilde{\psi}_{#1}^{\,#2}}

\begin{document}
\maketitle

\section*{Setup}

Symmetry acts on physical states; the relation is between
all-electron (AE) orbitals. For a point-group operation $\hat{R}$
mapping IBZ point $k_1$ to full-BZ point $k_2 = R\,k_1$,
\begin{equation}
  \psiAE{n}{k_2}(\mathbf{r})
  \;=\; \sum_{m} D_{mn}(R, k_1)\,
        (\hat{R}\,\psiAE{m}{k_1})(\mathbf{r}),
  \qquad
  (\hat{R}\,\psiAE{m}{k_1})(\mathbf{r})
  = \psiAE{m}{k_1}(R^{-1}\mathbf{r}).
  \label{eq:basis-relation-AE}
\end{equation}
Define
\begin{align}
  F_{pn}(R, k_1) &\;:=\; \matel{\rotpsiAE{p}{k_1}}{}{\psiAE{n}{k_2}}_{\AEt},
  \\
  M_{pm}(R, k_1) &\;:=\; \matel{\rotpsiAE{p}{k_1}}{}{\rotpsiAE{m}{k_1}}_{\AEt},
\end{align}
where $\langle\cdot|\cdot\rangle_{\AEt}$ is the inner product on the
all-electron Hilbert space. Combining with
Eq.~\eqref{eq:basis-relation-AE},
\begin{equation}
  F \;=\; M\,D, \qquad
  D \;=\; M^{-1}\,F.
  \label{eq:master}
\end{equation}

In NCPP the AE orbitals coincide with the pseudo ones,
$\psiAE{n}{k} \equiv \psipt{n}{k}$, and the AE inner product reduces
to the plane-wave dot product on the smooth coefficients. In USPP/PAW
the link between the AE and the pseudo orbital is the
projector-augmentation transformation
$\ket{\psiAE{n}{k}} = \widehat{T}\,\ket{\psipt{n}{k}}$, with
\begin{equation}
  \widehat{T}
  \;=\; 1 \;+\; \sum_{a,L}
        \big(\ket{\phi_L^{\AEt,a}} - \ket{\phi_L^{\mathrm{PS},a}}\big)
        \bra{p_L^{a}},
\end{equation}
so that the AE inner product becomes the augmented pseudo inner
product
\begin{equation}
  \matel{\psiAE{a}{k}}{}{\psiAE{b}{k}}_{\AEt}
  \;=\; \matel{\psipt{a}{k}}{S^{(k)}}{\psipt{b}{k}},
  \qquad
  S^{(k)} \;=\; \widehat{T}^{\dagger}\widehat{T}
  \;=\; 1 + \sum_{a,L,L'} \ket{p_L^{a,k}}\,Q^a_{LL'}\,\bra{p_{L'}^{a,k}},
  \label{eq:S-aug}
\end{equation}
with $Q^a_{LL'} = \matel{\phi_L^{\AEt,a}}{}{\phi_{L'}^{\AEt,a}}
- \matel{\phi_L^{\mathrm{PS},a}}{}{\phi_{L'}^{\mathrm{PS},a}}$.

\section*{Case 1 --- NCPP, orthonormal AE bands}

$\widehat{T} = 1$, $\psiAE{n}{k} \equiv \psipt{n}{k}$,
$\matel{\psiAE{p}{k}}{}{\psiAE{q}{k}}_{\AEt} = \delta_{pq}$. Then
$M = \mathbb{I}$ and
\begin{equation}
  \boxed{
  D_{pn}(R, k_1)
  \;=\; \sum_{\mathbf{G}}
        \big[\rotpsipt{p}{k_1}(\mathbf{G})\big]^{*}\,
        \psipt{n}{k_2}(\mathbf{G}).
  }
  \label{eq:D-orth}
\end{equation}

\section*{Case 2 --- USPP/PAW, AE-orthonormal bands}

The QE-shipped bands satisfy
$\matel{\psiAE{p}{k}}{}{\psiAE{q}{k}}_{\AEt}
= \matel{\psipt{p}{k}}{S^{(k)}}{\psipt{q}{k}}
= \delta_{pq}$, so $M = \mathbb{I}$ and
\begin{equation}
  \boxed{
  \begin{aligned}
  D_{pn}(R, k_1)
  \;=\; & \matel{\rotpsiAE{p}{k_1}}{}{\psiAE{n}{k_2}}_{\AEt}
  \;=\; \matel{\rotpsipt{p}{k_1}}{S^{(k_2)}}{\psipt{n}{k_2}}
  \\[2pt]
  \;=\; & \sum_{\mathbf{G}}
        \big[\rotpsipt{p}{k_1}(\mathbf{G})\big]^{*}\,
        \psipt{n}{k_2}(\mathbf{G})
  \\[2pt]
  & +\;
  \sum_{a, L, L'}
        \matel{\rotpsipt{p}{k_1}}{}{p_L^{a, k_2}}\,
        Q^a_{LL'}\,
        \matel{p_{L'}^{a, k_2}}{}{\psipt{n}{k_2}}.
  \end{aligned}
  }
  \label{eq:D-paw}
\end{equation}

\section*{Case 3 --- non-orthogonal AE bands}

Drop the AE-orthonormality assumption; keep both $F$ and $M$ in the
master equation~\eqref{eq:master}:
\begin{equation}
  \boxed{
  D(R, k_1) \;=\; \big[M(R, k_1)\big]^{-1}\,F(R, k_1).
  }
  \label{eq:D-nonorth}
\end{equation}
Written in the pseudo basis (covers NCPP, USPP, and PAW
simultaneously),
\begin{align}
  F_{pn} &\;=\;
        \matel{\rotpsipt{p}{k_1}}{S^{(k_2)}}{\psipt{n}{k_2}}
        \notag \\
        &\;=\;
        \sum_{\mathbf{G}} \big[\rotpsipt{p}{k_1}(\mathbf{G})\big]^{*}
                          \psipt{n}{k_2}(\mathbf{G})
        \;+\;
        \sum_{a, L, L'}
              \matel{\rotpsipt{p}{k_1}}{}{p_L^{a, k_2}}\,
              Q^a_{LL'}\,
              \matel{p_{L'}^{a, k_2}}{}{\psipt{n}{k_2}},
  \\[4pt]
  M_{pm} &\;=\;
        \matel{\rotpsipt{p}{k_1}}{S^{(k_2)}}{\rotpsipt{m}{k_1}}
        \notag \\
        &\;=\;
        \sum_{\mathbf{G}} \big[\rotpsipt{p}{k_1}(\mathbf{G})\big]^{*}
                          \rotpsipt{m}{k_1}(\mathbf{G})
        \;+\;
        \sum_{a, L, L'}
              \matel{\rotpsipt{p}{k_1}}{}{p_L^{a, k_2}}\,
              Q^a_{LL'}\,
              \matel{p_{L'}^{a, k_2}}{}{\rotpsipt{m}{k_1}}.
\end{align}
For NCPP the augmentation terms vanish ($Q \equiv 0$); the equations
reduce to the smooth dot products of the pseudo (= AE) coefficients.

\section*{Evaluating the augmentation term}

\paragraph{Should this involve $\widehat{T}_{k_1}^{\dagger}\,\widehat{T}_{k_2}$?}
No --- $\widehat{T}$ is a single, real-space, $k$-independent operator,
$\widehat{T} = 1 + \sum_{a,L}(\ket{\phi_L^{\AEt,a}} -
\ket{\phi_L^{\mathrm{PS},a}})\bra{p_L^{a}}$, where the sums run over
all atoms in the crystal (every periodic image) and the partial
waves/projectors are the same atom-localised functions used for every
$k$. The apparent $k$-dependence enters only through Bloch summation
over the periodic images when one projects a Bloch state against
those atom-local projectors.

In our case both states in
$F_{pn} = \matel{\rotpsipt{p}{k_1}}{\widehat{T}^{\dagger}\widehat{T}}{\psipt{n}{k_2}}$
carry crystal momentum $k_2$: $\rotpsipt{p}{k_1}$ has crystal
momentum $R\,k_1 = k_2$ by construction, and $\psipt{n}{k_2}$ is at
$k_2$ already. Performing the sum over all atom positions
$\tau_\alpha + T$ with $\alpha$ a unit-cell atom and $T$ a lattice
translation, the Bloch translation phases combine as
$e^{-i k_2 \cdot T} e^{+i k_2 \cdot T} = 1$, leaving
\begin{equation}
  F_{pn}^{\mathrm{aug}}
  \;=\; \sum_{\alpha, L, L'}
        \matel{\rotpsipt{p}{k_1}}{}{p_L^{\alpha, k_2}}\,
        Q^{\alpha}_{LL'}\,
        \matel{p_{L'}^{\alpha, k_2}}{}{\psipt{n}{k_2}},
  \label{eq:F-aug-explicit}
\end{equation}
where the sum is over unit-cell atoms $\alpha$ and the projectors
$\ket{p_L^{\alpha, k_2}}$ are the standard Bloch-summed,
unit-cell-restricted projectors at $k_2$:
\begin{equation}
  \ket{p_L^{\alpha, k}}
  \;=\; \frac{1}{\sqrt{N_{\mathrm{cell}}}}
        \sum_{T} e^{i k \cdot T}\,
        \ket{p_L^{\alpha}(\mathbf{r}-\boldsymbol{\tau}_\alpha-\mathbf{T})}.
\end{equation}
The right-hand projection
$\matel{p_{L'}^{\alpha, k_2}}{}{\psipt{n}{k_2}}$ is the standard
\texttt{Pskna} entry at $k_2$, already in storage. The new piece is
the left-hand projection
$\matel{\rotpsipt{p}{k_1}}{}{p_L^{\alpha, k_2}}$, between an orbital
indexed by $k_1$ (its storage label) and a projector at $k_2$ (its
crystal-momentum label after rotation).

\paragraph{Why $k_1 \neq k_2$ does not break the calculation.}
If we did the same exercise with $\matel{\psipt{p}{k_1}}{S}{\psipt{n}{k_2}}$
\emph{without} the rotation, the Bloch-translation phase sum gives
$\sum_T e^{i (k_2 - k_1)\cdot T} = N_{\mathrm{cell}}\,\delta_{k_1, k_2 \bmod \mathbf{G}}$,
which vanishes for distinct $k_1, k_2$ as it must for an orthonormal
Bloch basis. The rotation $\hat{R}$ shifts the bra-side crystal
momentum from $k_1$ to $k_2$, restoring the matching that makes the
overlap non-trivial.

\paragraph{Strategy A --- direct re-projection at $k_2$.}
Apply the same $\mathbf{G}$-vector rotation that builds
$\rotpsipt{p}{k_1}(\mathbf{G})$ for the smooth term (already in the
code at \texttt{transform\_k2g}), and reuse the standard projector
machinery at $k_2$ to compute
\begin{equation}
  \matel{\rotpsipt{p}{k_1}}{}{p_L^{\alpha, k_2}}
  \;=\; \sum_{\mathbf{G}} \big[\rotpsipt{p}{k_1}(\mathbf{G})\big]^{*}\,
        p_L^{\alpha, k_2}(\mathbf{G})\, e^{-i (\mathbf{G}+k_2)\cdot \boldsymbol{\tau}_\alpha}.
\end{equation}
Pros: reuses the FFT-rotated coefficients already needed for the
smooth term; no symmetry-property machinery for the projectors. Cons:
costs an additional projector sweep per $(R, k_1)$ pair.

\paragraph{Strategy B --- symmetry-routed projection at $k_1$.}
Use the transformation property of the projector under $\hat{R}$. With
$R$ acting on atom $\alpha$ as
$\hat{R}\,(\boldsymbol{\tau}_\alpha) = \boldsymbol{\tau}_{R\alpha}
+ \mathbf{t}_R^{\alpha}$ (the residual translation $\mathbf{t}_R^{\alpha}$
takes the rotated atom into the same unit cell), and $D^{(l)}_{ML}(R)$
the Wigner $D$-matrix of the angular momentum $l = l_L$:
\begin{equation}
  \hat{R}\,\ket{p_L^{\alpha, k_1}}
  \;=\; e^{-i k_2 \cdot \mathbf{t}_R^{\alpha}}
        \sum_{M} D^{(l)}_{ML}(R)\,
        \ket{p_M^{R\alpha,\, k_2}}.
\end{equation}
Hence
\begin{equation}
  \matel{\rotpsipt{p}{k_1}}{}{p_L^{\alpha, k_2}}
  \;=\; \matel{\psipt{p}{k_1}}{\hat{R}^{\dagger}}{p_L^{\alpha, k_2}}
  \;=\; e^{+i k_2 \cdot \mathbf{t}_R^{R^{-1}\alpha}}
        \sum_{M}\!\big[D^{(l)}_{ML}(R)\big]^{*}\,
        \matel{\psipt{p}{k_1}}{}{p_M^{R^{-1}\alpha,\, k_1}}.
  \label{eq:strat-B}
\end{equation}
Pros: the projector overlaps on the right of Eq.~\eqref{eq:strat-B}
are exactly the standard \texttt{Pskna} entries at $k_1$, already
cached --- no FFT rotation of the wavefunction, no recomputation of
projector convolutions. Cost is only the per-$(R,\alpha,L)$ Wigner
rotation. Cons: requires the lookup tables $R\alpha$,
$\mathbf{t}_R^{\alpha}$, $D^{(l)}(R)$ to be in place.

\paragraph{Strategy C --- precompute the rotated \texttt{Pskna} table.}
A practical hybrid: Strategy~B is applied once per $(R, k_1)$ pair to
build a rotated projector-overlap table
\begin{equation}
  \widetilde{P}^{(R, k_1)}_{p, L, \alpha}
  \;:=\; \matel{\rotpsipt{p}{k_1}}{}{p_L^{\alpha, k_2}}
\end{equation}
which is then reused for all $n$ in
Eq.~\eqref{eq:F-aug-explicit}. Cost: same as Strategy~B for the table
build, then a single \texttt{GEMM} per $(R, k_1)$ block to contract
with $Q^{\alpha}_{LL'}$ and the $k_2$-side
$\matel{p_{L'}^{\alpha, k_2}}{}{\psipt{n}{k_2}}$. This is the
recommended layout: it keeps the full $D$-matrix as a single GEMM
plus the augmentation correction as a second GEMM of the same shape.

\paragraph{Recap of what is needed.} Implementation footprint for the
QE-source PAW workflow ($M = \mathbb{I}$):
\begin{itemize}
  \item The species-resolved $Q^{\alpha}_{LL'}$ is already loaded in
        \texttt{pseudopot::qq\_nt\_data}.
  \item The right-hand $\matel{p_{L'}^{\alpha, k_2}}{}{\psipt{n}{k_2}}$
        is the standard \texttt{Pskna} entry at $k_2$ (already cached).
  \item The left-hand $\matel{\rotpsipt{p}{k_1}}{}{p_L^{\alpha, k_2}}$
        is the new piece. Strategy~A and~C share the same final
        contraction; they differ only in how the table is built (FFT
        rotation vs. Wigner-$D$ symmetry routing). Strategy~B/C is
        cheaper if the symmetry tables are already on hand.
\end{itemize}

\end{document}
